\chapter{Introduction}
\label{chap:intro}

\section{Motivation}
\label{sec:intro:motivation}

Realizing a Quantum Internet~\cite{VanMeter2014Quantum,Kozlowski2023Architectural}, i.e. a network connecting quantum devices over long distances, would enable applications with no classical equivalent~\cite{Azuma2023Quantum,Hajdusek2023Quantum}. \Gls{qkd} provides a shared secret key whose security is guaranteed unconditionally by the principles of quantum mechanics, independently of an eavesdropper's computational capabilities~\cite{Bennett1992Quantum,Bennett2014Quantum,Scarani2009Security}. \Gls{dqc} lets several small quantum processors connected by such a network act together as one larger one~\cite{Cirac1999Distributed,Boschero2025Distributed}. Finally, networked quantum sensing can improve measurement precision beyond what independent classically connected sensors can achieve~\cite{Azuma2023Quantum,Proctor2017Networked,Proctor2018Multiparameter}. These applications rely on distributed entanglement as a shared quantum resource; in the repeater architectures considered here, this resource is represented by high-fidelity Bell pairs, also called \emph{ebits}. Direct photon transmission cannot supply this resource over long distances because channel loss grows exponentially with length~\cite{Azuma2023Quantum}. Quantum repeaters address this limitation~\cite{Briegel1998Quantum}; Chapter~\ref{chap:relatedwork} reviews how successive repeater generations do so.

Regardless of repeater generation, the delivered entanglement is noisy. Entanglement purification is one standard way to turn such noisy entanglement into a usable resource: it consumes several independently generated copies of a shared pair and returns fewer copies of higher fidelity~\cite{Bennett1996Purification,Deutsch1996Quantum}. Each round trades additional resource consumption for a probabilistic fidelity gain. In a multi-hop repeater chain, purification is not restricted to a single point in the protocol: its placement and number of rounds can be chosen as part of the overall protocol design. The number of copies consumed and the purification circuit may also vary from one stage to another. This thesis studies an all-photonic repeater architecture built from \Glspl{rgs}~\cite{Azuma2015All} (Chapter~\ref{chap:background} introduces its physical building blocks), where purification may be applied during resource generation, after individual links are established, or after end-to-end entanglement is created. Choosing how to allocate a fixed purification budget across this chain is therefore a combinatorial scheduling problem. The central question is which schedule gives the best specified trade-off between rate, fidelity, and resource cost.

\section{The Scheduling Gap}
\label{sec:intro:gap}

Two recent papers by Benchasattabuse, Hajdušek, and Van Meter establish the specific architecture and scheduling flexibility this thesis builds on. Throughout this thesis,~\cite{Benchasattabuse2026Bridging} is called \emph{the Bridging paper} and~\cite{Benchasattabuse2025Integrating} is called \emph{the Integrating paper}. Together, they are called \emph{the source papers}.

The Bridging paper introduces the \Gls{hrgs-long} building block: half of an \Gls{rgs}, generated and held at a single station rather than requiring the full state to exist in one place, allowing an end node to act as its own source of entanglement. This enables the all-photonic scheme to interoperate with conventional memory-equipped repeater nodes instead of requiring every intermediate station to be purely photonic~\cite{Benchasattabuse2026Bridging}. The Integrating paper then shows that entanglement purification can be inserted into this architecture at any stage, generation, link, or end node, and mixed freely between stages, provided it is executed \emph{optimistically}~\cite{Benchasattabuse2025Integrating,Mobayenjarihani2023Optimistic}, i.e. without waiting for a classical heralding outcome between rounds. Chapter~\ref{chap:background} formalises why this timing model is what makes purification worth inserting at all.

The Integrating paper demonstrates this flexibility with a single chosen schedule and identifies the search for a better one, under a fixed physical network and resource budget, as an open problem. It does not establish which allocation of a fixed budget performs best, or whether that answer changes with the noise model or chain length. This is the scheduling gap this thesis addresses. Given a fixed network configuration, the research problem is to select a schedule from the combinatorially many possibilities, which may differ in purification location, copy count, circuit, and order, and optimise it for a specified objective such as rate, fidelity, or cost. The working hypothesis is that systematic search finds schedules that improve on the hand-picked baseline at equal cost and, more importantly, retain feasibility in regimes where that baseline fails. This thesis tests that hypothesis across several network configurations beyond the source papers' reference configuration. Chapters~\ref{chap:model} and~\ref{chap:algorithms} formalise this schedule space and the algorithms that search it, respectively.

\section{Research Objectives and Contributions}
\label{sec:intro:contributions}

Addressing this gap first requires a precise object to search over, since neither source paper formalises the schedule space. Second, it requires search algorithms capable of exploring that exponentially large space. This thesis makes four concrete contributions toward that end.

\begin{enumerate}
    \item A formal, \Gls{dag}-based model of the schedule space for \Gls{hrgs-long}-based quantum repeaters (Chapter~\ref{chap:model}) and a bottom-up evaluator that reproduces the source papers' fidelity benchmarks for its raw, heralded, and optimistic end-node pumping schedules (defined in Chapter~\ref{chap:algorithms}) to four significant figures.

    \item A three-tier search framework: structured brute-force baselines, an exact span-partition dynamic program, and bounded same-span pumping and beam-search extensions (Chapter~\ref{chap:algorithms}). Every candidate schedule is structurally validated and evaluated using the same physical model. At the Integrating paper's reference configuration, the matched-cost result improves rate by $2.5\%$, while a budget-relaxed result reaches a rate of $6195.95$ at cost $60$, a $52.8\%$ improvement while only using $60\%$ of the paper baseline's budget (Chapter~\ref{chap:experiments}).

    \item An empirical study of when scheduling changes feasibility rather than only quality, varying outer- and inner-qubit noise, the number of \Gls{bsm} arms per hop, hop count, memory decoherence, and generation timing (Chapters~\ref{chap:experiments}--\ref{chap:discussion}). Its strongest result is a budget-fixed feasibility advantage: under non-ideal inner-qubit error, the paper schedule falls below the $F \ge 0.9$ threshold while a searched schedule at the same budget still clears it.

    \item A synthesis of when scheduling matters (Chapter~\ref{chap:discussion}): the advantage is largest when the budget is tight enough for purification choices to affect feasibility and rate, but not so restrictive that no schedule remains feasible. Within the searched families, the minimum feasible budget over $N=10$--$28$ grows approximately as $E_{\max}^{\min} \approx 1.19\,N^{1.67}$, suggesting a superlinear trend rather than the linear $10N$ budget assumed in the Integrating paper. This beam-search result is beam-width sensitive and should not be interpreted as a proven asymptotic bound.

\end{enumerate}

\noindent Every result above only establishes existence within the schedule families actually searched, rather than providing a global optimality guarantee; Section~\ref{sec:disc:scope} returns to this distinction.